\vspace{-0.15cm}
\section{Introduction} 
\vspace{-0.15cm}

Wind energy has demonstrated to be an excellent alternative energy source, thanks to being completely renewable and environmentally sustainable~\cite{windreport}. The increase in installed capacity urges to adapt to more demanding challenges in terms of cost and logistics~\cite{30capacity}, with special focus on developing non-destructive inspection manners for health monitoring. Drone inspections have become a fast and reliable setup for wind turbine diagnosis~\cite{drone}, capturing around 400 high-resolution images per turbine to design a repair plan. All in all results in high-volume data that encourages the development of automated, scalable and cost-effective solutions. 

To successfully design a data-driven maintenance framework for wind turbines in operation, current researchers are tackling distinct image-based problems such as blade defect detection~\cite{aidrone,maskrcnn} or autonomous navigation~\cite{realtime, automateddrone}. However, current solutions have several limitations and are not being applied yet to the industry. Hence, these challenges must be simplified by means of distinct low-level vision techniques. To this aim, this paper proposes addressing segmentation to reduce the complexity of these learning tasks.

%Segmenter: Transformer for Semantic Segmentation
%Image Segmentation Using Deep Learning: A Survey
\begin{figure}[t!]
%\begin{minipage}[b]{0.94\linewidth}
  \centering
\resizebox{8.4 cm}{!} {
  \centerline{\includegraphics[width=8.5cm]{Figures/introd.png}}}
  \vspace{-0.3cm}
%\end{minipage}
\caption{\textbf{Wind turbine blade images in the wild} with distinct complexity due to the variation of the landscapes, the blade shape and size, and the presence of shadows or dirt.}
\label{fig:introd}
\vspace{-0.55cm}
\end{figure}

Current state-of-the-art segmentation methods rely on learning-based approaches, usually employing an encoder-decoder architecture~\cite{deeplab,sw,cvpr2022} and, recently, an attention mechanism on the encoder~\cite{bmvc,attention2}. With a properly tailored architecture, these models have proven to successfully obtain high precision results. Unfortunately, they do not provide a segmentation algorithm that generalizes for newly acquired images~\cite{generic}. A generic solution is essential to deal with a large variety of images (see Fig.~\ref{fig:introd}) obtained on regular wind turbine inspections under a different setup. 

A common technique to enhance the segmentation mapping is including a post-processing step. DeepLab~\cite{deeplab} utilized Conditional Random Fields to fine-tune the localized segmented objects. DeepLab V3+~\cite{deeplabv3+} is further strengthened through superpixel merging in~\cite{generic}. Another interesting approach is ResNeSt~\cite{resnest}, which last stage incorporates a channel-wise attention with multi-path representation. However, these methods do not leverage our prior information about blade images, such as that the blade traverses the image or that images of the same inspection share image properties.  

This work proposes a novel segmentation algorithm that identifies the blade region for very diverse images in appearance. First, a preliminary mask is obtained by a customized U-Net model~\cite{unet}, whose architecture still inspires current techniques~\cite{unetformer}. Then, this preliminary result is improved through a set of post-processing steps: an extended hole filling routine from~\cite{fill} followed by an unsupervised random forest~\cite{randomforest}; refined again through another hole filling step. The hole filling relies on properly identifying the blade image borders and the random forest on underfitting the preliminary masks from the same blade surface. As a result, we obtain an accurate, robust and efficient model for blade segmentation that can handle a wide variety of images in the wild.
